\documentclass{anvil-paper}
\usepackage{array}
\usepackage{tabularx}
\usepackage{fontspec}
\setmainfont{Times New Roman}
\setmonofont{Menlo}[Scale=0.85]
\usepackage[htt]{hyphenat}

% Lean identifiers: typewriter with url-style breaking at underscores and dots,
% plus a discouraged (penalty 500) break before each capital letter so that long
% camel-case identifiers can break without a badness-10000 line.
% (No pandoc escapes: the argument is written with raw underscores.)
\makeatletter
\def\lean@camel{%
  \do\A{\penalty500\mathchar`A}\do\B{\penalty500\mathchar`B}\do\C{\penalty500\mathchar`C}%
  \do\D{\penalty500\mathchar`D}\do\E{\penalty500\mathchar`E}\do\F{\penalty500\mathchar`F}%
  \do\G{\penalty500\mathchar`G}\do\H{\penalty500\mathchar`H}\do\I{\penalty500\mathchar`I}%
  \do\J{\penalty500\mathchar`J}\do\K{\penalty500\mathchar`K}\do\L{\penalty500\mathchar`L}%
  \do\M{\penalty500\mathchar`M}\do\N{\penalty500\mathchar`N}\do\O{\penalty500\mathchar`O}%
  \do\P{\penalty500\mathchar`P}\do\Q{\penalty500\mathchar`Q}\do\R{\penalty500\mathchar`R}%
  \do\S{\penalty500\mathchar`S}\do\T{\penalty500\mathchar`T}\do\U{\penalty500\mathchar`U}%
  \do\V{\penalty500\mathchar`V}\do\W{\penalty500\mathchar`W}\do\X{\penalty500\mathchar`X}%
  \do\Y{\penalty500\mathchar`Y}\do\Z{\penalty500\mathchar`Z}}
\DeclareUrlCommand\lean{\urlstyle{tt}%
  \expandafter\def\expandafter\UrlSpecials\expandafter{\UrlSpecials\lean@camel}}
% Table variant: breaks only at underscores and dots, never inside a camel-case
% name, so that identifiers in narrow table cells stay copyable.
\DeclareUrlCommand\leantab{\urlstyle{tt}}
\makeatother
% Public links to the branch of record (BRIEF rule R-LINK). \repobase is the one
% base URL and is switched from the branch to the release tag at publication;
% \repofile{<path>}{<label>} links <label> (typewriter via \lean: url-style breaking
% at underscores, dots and slashes plus the discouraged camel-case breaks, so long
% module and file names can break) to \repobase/<path>; \repofiletab is the table
% variant (typewriter via \leantab, no camel-case breaks). hyperref is loaded by the class.
\newcommand{\repobase}{https://github.com/rjwalters/lean-genius/blob/erdos85/integration}
\newcommand{\repofile}[2]{\href{\repobase/#1}{\lean{#2}}}
\newcommand{\repofiletab}[2]{\href{\repobase/#1}{\leantab{#2}}}
% Long identifiers: allow moderate stretch instead of \sloppy.
\tolerance=800
\setlength{\emergencystretch}{3em}
% Keep tables with text rather than on float pages.
\renewcommand{\topfraction}{.9}
\renewcommand{\bottomfraction}{.9}
\renewcommand{\textfraction}{.05}
\renewcommand{\floatpagefraction}{.85}
\newcommand{\AREG}{\mathrm{A\text{-}REG}}
\newcommand{\NONBIP}{\mathrm{A\text{-}REG\text{-}NONBIP}}
\newcommand{\Qn}{\textsf{Erdos85Question}}
\newcommand{\hnoFortyNine}{\lean{hno49}}

\newtheorem*{resultA}{Result A}
\newtheorem*{theoremB}{Theorem B}
\theoremstyle{definition}
\newtheorem*{hypAREG}{Hypothesis A-REG}
\newtheorem*{nonbipconn}{NONBIP-CONNECTED}

\title{Computational Evidence for a Drop and a Uniform Reduction for Erdős Problem 85}
\author{Claude Fable and GPT Sol \\ Lean Genius project, 2AM Logic}
\date{28 September 2026}

\begin{document}
\maketitle

\begin{abstract}
For $n\ge 4$ let $f(n)$ be the least $d$ such that every simple graph on $n$ vertices with minimum degree at least $d$ contains a four-cycle; Erdős Problem 85 asks whether $f$ is eventually nondecreasing. We report two results. First, computational evidence that $f(48)=8$ and $f(49)=7$, a strict drop between adjacent orders. The lower sides are explicit graphs checked in Lean 4. The upper side at order 49 is a case split, proved exhaustive in Lean, by the number of degree-8 vertices: three strata are closed by checked proof certificates and reviewed arguments with independently checked computation, the fourth by a census in which two independent SAT solvers refuted every instance under declared caps, the hardest through an exact partition into 36 cubes, with no satisfiable instance and no disagreement; every one of that stratum's 1,257 instances then received an LRAT proof accepted by a formally verified checker. This is computational evidence, not a theorem: the checks run outside Lean and parts of the formal assembly remain open, and we state the evidence level of every input. To our knowledge it is the first strict drop of $f$ settled, at the level of computational evidence, on the problem's stated domain; in Boza's Ramsey convention it corresponds to selecting $r(42)=49$ in the open entry $r(42)\in\{49,50\}$. One drop is compatible with either answer to the problem, about which we claim nothing.

Second, Lean 4 verifies from the standard axioms alone that one uniform proposition, $\AREG$ (no $C_4$-free $2^k$-regular graph on $4^k$ vertices for any $k\ge 3$), implies a negative answer to Erdős Problem 85. $\AREG$ is an unproved hypothesis with a stated rival, not a forecast: the Lean-checked values $f(15)=f(16)=5$ refute its $q=4$ analogue.
\end{abstract}

\section{Introduction}\label{sec:intro}

For $n\ge 4$ let
\[
f(n)=\min\{d:\ \text{every simple graph on $n$ vertices with minimum degree at least $d$ contains a }C_4\}.
\]
Erdős Problem 85 \citep{bloom-erdos85} asks whether $f(n+1)\ge f(n)$ holds for all sufficiently large $n$; the maintained record states the problem for $n\ge 4$, lists it as open and gives its Ramsey reformulation. In Lean the function is \lean{minDegreeForC4} and the question is the proposition \Qn. The formal negation is equivalent to the existence of arbitrarily late strict drops (\lean{erdos85Negation_iff_not_question}), so a negative answer needs infinitely many drops, and a single drop decides nothing about the problem as posed.

The Ramsey side uses $r(s)=R(C_4,K_{1,s})$ \citep{boza2024ramsey,zhang2017polarity}. The two functions are related by an elementary conversion: $r(s)\le N$ holds exactly when every $C_4$-free graph on $N$ vertices has a vertex of degree at most $N-1-s$, that is, when $f(N)\le N-s$; hence $r(s)=\min\{N: f(N)\le N-s\}$, and a plateau $r(s)=r(s+1)=N$ says exactly that $f(N)\le N-s-1<N-s\le f(N-1)$, a strict drop from order $N-1$ to order $N$. We use $f$ throughout and say which convention is meant whenever a value from Boza's table is quoted.

\paragraph{Result A (computational).} Lean checks, through \texttt{native\_decide} on explicit graphs, that $f(48)=8$ (\lean{minDegreeForC4_fortyEight_eq_eight_checked}) and $f(49)\ge 7$ (\lean{seven_le_minDegreeForC4_fortyNine_checked}). The one unproved hypothesis is the nonexistence of a $C_4$-free graph on 49 vertices with minimum degree 7, the Lean hypothesis
\[
\text{\hnoFortyNine} \;:\; \neg\,\text{\texttt{C4FreeMinDegreeWitness 49 7}} .
\]
Under that hypothesis the Lean-checked finite-drop core gives $f(49)=7<8=f(48)$. Result A is the claim that a completed two-solver census supports the hypothesis at the level of computational evidence. In a hypothetical $C_4$-free graph on 49 vertices with minimum degree 7 every degree is 7 or 8, and Lean proves that the number $h$ of degree-8 vertices is odd and at most 9 and that $h=9$ is impossible (Section~\ref{sec:strata}); the upper side therefore splits into the strata H1, H3, H5 and H7. H3 is closed by checked LRAT proofs of its two cell formulas and, independently, by a reviewed argument; H5 by reviewed reductions with independently checked finite computation; H7 by thirteen LRAT certificates checked inside Lean by \texttt{native\_decide}, one for each of its thirteen canonical representatives with a positive triple, together with a reviewed structural closure of its last cell (Section~\ref{sec:upper}). H1 is 1,257 SAT instances emitted by a Lean program: 96 carry 2026-08 drat-trim-verified certificates, and 1,161 were refuted, without proof logging, by Kissat 4.0.4 and then CaDiCaL 3.0.1 under declared caps (1,160 as whole instances and one, which defeated every whole-instance cap up to 24 hours, as an exact partition into 36 cubes). The census recorded 1,412 input preparations over the four cloud passes for the 1,137 rows dispatched to the cloud (the 24 pilot rows ran on the local Mac; a row that reached its cap was prepared again for the next pass) and spent about 5,700 core-hours of final-attempt solver time; no row is satisfiable, no row is open, and the solvers never disagreed. In October 2026 every H1 instance was then certificate-checked: each of the 1,160 whole-instance rows was re-solved by CaDiCaL with proof logging and its binary LRAT proof checked, as it streamed, by \texttt{cake\_lpr}, a checker verified in CakeML \citep{tan2021cakelpr}; the 36 cube leaves of the hardest row were checked the same way and composed by a Lean lemma; and the 96 historical certificates were re-checked (Section~\ref{sec:cost}). Result A is evidence, not a theorem: Section~\ref{sec:trust} states where the residual trust sits.

\paragraph{Significance.} To our knowledge, this is the first strict drop of $f$ settled, at the level of computational evidence, on the problem's stated domain $n\ge 4$: the maintained record \citep{bloom-erdos85} lists no partial solution and warns that its literature list may be incomplete, and Boza's decided table \citep{boza2024ramsey} reports no consecutive equality that yields an in-domain drop, though some of its entries remain ranges, so an earlier undecided candidate cannot be ruled out. In Boza's convention the relevant entries are $r(41)=49$ and $r(42)\in\{49,50\}$; the nonexistence hypothesis selects $r(42)=49$, giving the plateau $r(41)=r(42)=49$, which by the conversion above is exactly $f(49)=7<8=f(48)$. A satisfying assignment in any H1 row would instead have selected $r(42)=50$ and there would be no drop at 49. One drop is compatible with either answer to Erdős Problem 85, and we make no claim about the problem itself.

\paragraph{Theorem B (Lean, standard axioms only).} Let $\AREG$ (\lean{BinarySquareRegularExclusion}) assert that for every $k\ge 3$ there is no simple $C_4$-free $2^k$-regular graph on $4^k$ vertices. Lean 4 \citep{moura2021lean4} with mathlib \citep{mathlib2020} verifies
\[
\AREG \;\Rightarrow\; \neg\,\Qn
\]
as \lean{not_erdos85Question_of_binarySquareRegularExclusion}, depending on no axioms beyond \lean{propext}, \lean{Classical.choice} and \lean{Quot.sound}. $\AREG$ is an unproved hypothesis with a stated rival, not a conjecture we endorse: Section~\ref{sec:theoremB} states what has been proved beneath it, what residue remains, and why the evidence is mixed.

\paragraph{What the paper does not claim.} Result A is never called a theorem, a proof or a decided value; no solver verdict or archived certificate is promoted to a kernel-checked statement; nothing is claimed about Erdős Problem 85 itself; and $\AREG$ is not asserted. Every number in Sections~\ref{sec:intro} to~\ref{sec:contrib} traces to a receipt file named in Section~\ref{sec:artifacts}, except that the plane-order census figures quoted in Section~\ref{sec:areg-evidence} trace to the ledger files linked beside them in Appendix~\ref{app:negative}. Every file, directory and Lean module the paper names is linked to its location on the branch of record (Section~\ref{sec:artifacts}), and material that is not published is described as such.

\paragraph{Organization.} Section~\ref{sec:related} places the work against Boza's table, the Afzaly--McKay records and the SAT-with-certificates lineage. Section~\ref{sec:prelim} fixes the Lean definitions, the strata and the formal interface. Section~\ref{sec:resultA} presents Result A: the witnesses, the closure of each stratum, the evidence-level table, where the residual trust sits, and the cost of a formal closure including the cube-partitioned route. Section~\ref{sec:theoremB} presents Theorem B and its residue. Section~\ref{sec:interpretation} interprets both results. Appendix~\ref{app:negative} condenses the negative map for $\AREG$; Appendix~\ref{app:collab} records the technical lessons of the collaboration.

\section{Related work}\label{sec:related}

\paragraph{The problem and its Ramsey form.} The maintained Erdős Problems record \citep{bloom-erdos85} states Problem 85 for $n\ge 4$, marks it open and records no partial solution. The star-Ramsey function $r(s)=R(C_4,K_{1,s})$ is the object of \citet{zhang2017polarity}, who relate it to polarity graphs, and of \citet{boza2024ramsey}, whose table gives exact values and bounds; the conversion in Section~\ref{sec:intro} turns a plateau of $r$ into a drop of $f$. We cite Boza for the table and the open entry $r(42)\in\{49,50\}$, and Zhang, Chen and Cheng for the polarity-graph values used in Appendix~\ref{app:negative}. We do not compare the methods behind the decided entries of the $r$ table with the method used here.

\paragraph{Extremal records at order 49.} \citet{afzaly-mckay-extremal} list six $C_4$-free graphs on 49 vertices with 174 edges and minimum degree 6, labelled on their page as the best examples found rather than as an exhaustive computation. They corroborate $f(49)\ge 7$ and are cited here only as lower-bound examples; they do not exclude a different 49-vertex $C_4$-free graph of minimum degree 7, which is the obligation Result A addresses. Our own 48-vertex witness is not isomorphic to any graph of the same archive at order 48 (Section~\ref{sec:lower}).

\paragraph{SAT solving with and without certificates.} The census uses Kissat \citep{biere2024kissat} and CaDiCaL \citep{biere2020cadical} as independent solvers. The distinction the paper rests on, between a solver's UNSAT verdict and a checkable proof of it, is the standard one: DRAT proofs are checked by drat-trim \citep{wetzler2014drat}, and the LRAT format \citep{cruzfilipe2017lrat} adds the hints that let a small verified checker replay them. Cube-and-conquer \citep{heule2011cube} partitions a hard instance into cubes that are solved independently; the Boolean Pythagorean triples proof \citep{heule2016pythagorean} is the model for a cube-partitioned result whose per-cube proofs are checked separately. The historical H1 rows follow that model (certificates checked by drat-trim at production time); the residual rows do not, and Section~\ref{sec:cube} projects what following it would cost.

\paragraph{The trust root.} Graph semantics, the finite-drop core, the stratum interface and Theorem B are stated and checked in Lean 4 \citep{moura2021lean4} over mathlib \citep{mathlib2020}. The finite witnesses and the thirteen H7 certificate checks are checked through \texttt{native\_decide}, which adds compiler-generated axioms: we list them for the witnesses (Table~\ref{tab:axioms}) and state the dependency for the certificate checks (Section~\ref{sec:upper}); Theorem B uses the standard axioms only.

\section{Definitions and the formal interface}\label{sec:prelim}

\subsection{Lean definitions and audited statements}\label{sec:lean-defs}

\texttt{minDegreeForC4 n} is $f(n)$. \texttt{C4FreeMinDegreeWitness n d} asserts that a simple $C_4$-free graph on $n$ vertices with minimum degree at least $d$ exists. The module \repofile{proofs/Proofs/Erdos85FiniteDropWitnesses.lean}{Erdos85FiniteDropWitnesses.lean} supplies the two explicit graphs \lean{boza48_degreeSeven_witness} and \lean{orderFortyNine_degreeSix_witness}; the finite-drop core consists of the two statements
\begin{flushleft}\small
\lean{minDegreeForC4_fortyEight_fortyNine_exact_checked} and \lean{minDegreeForC4_fortyNine_lt_fortyEight_checked},
\end{flushleft}
\noindent both of which take \hnoFortyNine\ as a hypothesis. A cold rebuild on 2026-09-27 (fresh build volume, pinned Lean 4.31.0 image, mathlib from the standard cache) printed the literal \texttt{\#print axioms} output summarized in Table~\ref{tab:axioms}. Lean records each \texttt{native\_decide} use as a per-declaration axiom named \texttt{\ldots\_native.native\_decide.ax\_N}; the six such entries in Table~\ref{tab:axioms} are instances of the compiler-trust axiom \texttt{Lean.ofReduceBool}, and that single name is used below for the thirteen H7 certificate checks of Section~\ref{sec:upper}, whose axiom lists were not printed in the cold audit. Two conditional endpoints under the planned names
\begin{flushleft}\small
\lean{minDegreeForC4_fortyEight_fortyNine_exact_of_generatedSevenBaseCertificates} and \lean{minDegreeForC4_fortyNine_lt_fortyEight_of_generatedSevenBaseCertificates}
\end{flushleft}
\noindent are specified in a generated module that is not present in the source tree; they would prove the finite drop only after every input is supplied, the module lands, and a cold build and literal axiom audit pass.

\begin{table}[htbp]
\centering\small
\caption{Audited Lean statements and their axiom scope (cold rebuild, 2026-09-27). Every entry also depends on \texttt{propext}, \texttt{Classical.choice} and \texttt{Quot.sound}. The thirteen H7 certificate modules of Section~\ref{sec:upper} were not part of this rebuild and are not listed.}
\label{tab:axioms}
\begin{tabularx}{\textwidth}{@{}>{\raggedright\arraybackslash}p{.34\textwidth}>{\raggedright\arraybackslash}p{.26\textwidth}>{\raggedright\arraybackslash}X@{}}
\toprule
Lean statement & Content & Axioms beyond the standard three \\
\midrule
\leantab{not_erdos85Question_of_binarySquareRegularExclusion} & Theorem B: $\AREG\Rightarrow\neg\,\Qn$ & none \\
\leantab{minDegreeForC4_fortyEight_eq_eight_checked} & $f(48)=8$ & three \texttt{native\_decide} axioms (the 48-vertex graph, its codegree bound, its degree) \\
\leantab{seven_le_minDegreeForC4_fortyNine_checked} & $f(49)\ge 7$ & three \texttt{native\_decide} axioms (the 49-vertex graph, its codegree bound, its degree) \\
\leantab{minDegreeForC4_fortyEight_fortyNine_exact_checked} & $f(48)=8\wedge f(49)=7$, given \hnoFortyNine & the six above \\
\leantab{minDegreeForC4_fortyNine_lt_fortyEight_checked} & $f(49)<f(48)$, given \hnoFortyNine & the three 48-vertex axioms \\
\bottomrule
\end{tabularx}
\end{table}

\subsection{The strata and the consumer interface}\label{sec:strata}

Let $G$ be a simple $C_4$-free graph on 49 vertices with minimum degree at least 7. Because two vertices have at most one common neighbour, the length-two walks from a vertex $v$ that do not return to $v$ have distinct endpoints, so $\sum_{u\sim v}(\deg u-1)\le 48$; with every degree at least 7 this forces every degree to be 7 or 8, and a vertex of degree 8 to have only neighbours of degree 7. Call a vertex \emph{high} if its degree is exactly 8 (\lean{orderFortyNineHighVertices G} is the finite set of such vertices) and write $h$ for their number. The degree sum $7\cdot 49+h$ is even, so $h$ is odd, and a counting bound proved in Lean gives $h\le 9$: Lean proves $h\in\{1,3,5,7,9\}$ for every $C_4$-free graph on 49 vertices with all degrees at least 7 (\lean{orderFortyNine_card_high_eq_one_or_three_or_five_or_seven_or_nine}). The proposition \lean{OrderFortyNineStratumExcluded h} states that no such graph on \texttt{Fin 49} has exactly $h$ high vertices. The $h=9$ stratum is refuted in Lean with no external input (\lean{orderFortyNineStratumExcluded_nine}, through \lean{false_of_orderFortyNine_nine_high}), and
\[
\begin{aligned}
&\text{\lean{not_c4FreeMinDegreeWitness_fortyNine_seven_of_strata}}\ (h_1)\,(h_3)\,(h_5)\,(h_7)\\
&\qquad :\;\neg\,\text{\texttt{C4FreeMinDegreeWitness 49 7}}
\end{aligned}
\]
composes the four remaining exclusions. The case split is therefore the theorem's own interface: its only inputs are the four stratum exclusions of Table~\ref{tab:inputs}, and the evidence of Section~\ref{sec:resultA} is evidence for those four propositions. Each stratum is further cut into cells by the high--low incidence pattern (Section~\ref{sec:upper}); Table~\ref{tab:inputs} names the Lean statement that reduces each stratum to its cells.

\begin{table}[htbp]
\centering\small
\caption{The strata of the order-49 case split, the Lean statement that reduces each to its cells, and the evidence supplied for its exclusion.}
\label{tab:inputs}
\begin{tabularx}{\textwidth}{@{}>{\raggedright\arraybackslash}p{.04\textwidth}>{\raggedright\arraybackslash}p{.47\textwidth}>{\raggedright\arraybackslash}X@{}}
\toprule
$h$ & Lean reduction to cells & Evidence for the exclusion \\
\midrule
1 & \leantab{orderFortyNineStratumExcluded_one_of_pureFamilies}: five ordered one-high family profiles, refined by the capacity rows of the H1 index & two-solver census (Section~\ref{sec:upper}); the row-to-stratum assembly is an open formal obligation \\
3 & \leantab{orderFortyNineStratumExcluded_three_of_tripleCells}: cells $t=0,1$ & checked LRAT proofs of \leantab{orderFortyNineGeneratedCanonicalSatCnf} at \leantab{threeHighRepresentativeMasks} for both cells; independent paper-and-computation exclusion \\
5 & \leantab{orderFortyNineStratumExcluded_five_of_tripleCells}: cells $t=0,1,2$ (T0, T1, T2) & reviewed reductions of the three cells with independently checked finite computations (reviews 2037, 2032, 2062 and 2063; outer ledger review 2065, 2026-09-10); the three Boolean-exclusion premises of \leantab{orderFortyNineStratumExcluded_five_of_booleanExclusions} are undischarged in Lean \\
7 & \leantab{orderFortyNineStratumExcluded_seven_of_tripleCells}: cells $t=0,\dots,7$ (\leantab{orderFortyNine_highIncidence_profile_of_seven_high} bounds $t$ by 7); fourteen canonical representatives, $1,1,2,3,3,2,1,1$ by $t$, from a kernel-checked census and the proved graph cover \leantab{sevenHighCanonicalGraphCover_all} & cells $t=1,\dots,7$: thirteen LRAT certificates checked in Lean by \texttt{native\_decide} (2026-08-15 manifests), so that \leantab{orderFortyNineStratumExcluded_seven_of_t0} reduces the stratum to the $t=0$ representative; cell $t=0$: closure ledger of 2026-09-15; the Lean capstone for that cell is uninstantiated \\
9 & \leantab{orderFortyNineStratumExcluded_nine} & proved in Lean, no input \\
\bottomrule
\end{tabularx}
\end{table}

Two Lean consumers reach \hnoFortyNine\ from different H3 and H5 inputs, and their input counts must not be mixed or added.
\begin{itemize}
\item The LRAT consumer, \lean{not_c4FreeMinDegreeWitness_fortyNine_seven_of_smallHighLratChecks} in \repofile{proofs/Proofs/Erdos85OrderFortyNineSmallHighVerifiedFrontier.lean}{Erdos85OrderFortyNineSmallHighVerifiedFrontier.lean}, takes five whole-cell inputs, checked LRAT proofs of the Lean-generated formula \lean{orderFortyNineGeneratedCanonicalSatCnf} at the two H3 representative masks (\lean{threeHighRepresentativeMasks}, cells $t=0,1$) and at the three H5 representative masks (\lean{fiveHighRepresentativeMasks}, cells $t=0,1,2$), together with the H1 and H7 exclusions. Its docstring records that the H3 and H5 graph-normalization obligations, that every graph of a cell satisfies its formula, are discharged. This is an interface description, not a claim that five whole-cell certificates are available: two are (Section~\ref{sec:upper}).
\item The cube consumer, \lean{not_c4FreeMinDegreeWitness_fortyNine_seven_of_smallHighCubeBaseUnsat} in \repofile{proofs/Proofs/Erdos85OrderFortyNineSmallHighCubeGridTerminal.lean}{Erdos85OrderFortyNineSmallHighCubeGridTerminal.lean}, takes seven base formulas instead: \lean{Unsat} for the four H3 scout formulas (B1, C1, C2 and distance 2) and for the three H5 cell formulas, each supplied by the checked cube-grid theorem from that base formula's cover and leaf evidence. It reaches the same conclusion without the whole-cell LRAT arrays, and the finite-drop core accepts either route.
\end{itemize}
H1 rows still need the full row-to-stratum assembly. H7's Lean interface is described with the stratum in Section~\ref{sec:upper}: thirteen of its fourteen canonical representatives carry certificates checked inside Lean, and the capstone for the last one, \lean{orderFortyNineStratumExcluded_seven_of_emptyCubeEvidenceVectors}, takes four evidence-vector arguments of lengths 19, 15, 7 and 2 and is not yet instantiated.

\section{Result A: computational evidence for a 48-to-49 drop}\label{sec:resultA}

\begin{resultA}
The combined evidence supports $f(48)=8$ and $f(49)=7$, and consequently the strict drop $f(49)<f(48)$.
\end{resultA}

This is a computational result, not an unconditional Lean theorem. The census rule was that any H1 row reaching its declared solver cap would be printed as open; after escalating caps and one cube partition, no row is open, and every H1 instance now also carries an LRAT proof accepted by a verified checker (Section~\ref{sec:cost}). A formal theorem would still need the H1 semantic bridge and the other obligations of Section~\ref{sec:evidence}.

\subsection{Lower sides: the two witnesses}\label{sec:lower}

The two witnesses of Section~\ref{sec:lean-defs} are elaborated in Lean through \texttt{native\_decide}; the cold-rebuild audit lists three \texttt{native\_decide} axioms per witness (Table~\ref{tab:axioms}), so the witness side is not a standard-axiom-only kernel proof. Independent edge-list audits corroborate the graphs: the 48-vertex witness has 168 edges and all pair codegrees at most 1, and the 49-vertex witness passes the same codegree check. Separately, a NetworkX 3.6.1 isomorphism comparison (\repofile{research/problems/erdos-85-wip-01/sat49/verify_boza48_nonisomorphism.py}{sat49/verify_boza48_nonisomorphism.py}, rerun 2026-09-28 with its receipt \repofile{research/problems/erdos-85-wip-01/sat49/BOZA48_NONISOMORPHISM_RECEIPT_20260928.txt}{sat49/BOZA48_NONISOMORPHISM_RECEIPT_20260928.txt}) finds the 48-vertex witness non-isomorphic to all ten graphs of the Afzaly--McKay archive \citep{afzaly-mckay-extremal} at order 48 with 168 edges, only one of which is 7-regular, so it is a second realization; this comparison is reproducible computation, not part of the Lean-elaborated theorem.

\subsection{Upper side: closing the four strata}\label{sec:upper}

\paragraph{H3 and H5: cells and whole-cell formulas.} Put the $h$ high vertices first and write the adjacency matrix as $A=\left[\begin{smallmatrix}0&B\\B^{\mathsf T}&C\end{smallmatrix}\right]$. High vertices are pairwise non-adjacent, each pair of them has exactly one common neighbour, and a low vertex has at most three high neighbours (its \emph{support}), since its neighbourhood induces a matching and every high neighbour must be matched to a distinct low neighbour. Up to relabelling the high vertices, the high--low incidence pattern falls into finitely many cells indexed by the number $t$ of low vertices of support 3. For $h=3$ there are two incidence profiles, with $(25,18,3,0)$ and $(24,21,0,1)$ low vertices of support $0,1,2,3$, one for each of the cells $t=0,1$ of Table~\ref{tab:inputs} (the canonical cells of the Lean statement refine the profiles by fixing a representative mask); for $h=5$ the three cells T0, T1 and T2 have support profiles $(14,20,10,0)$, $(13,23,7,1)$ and $(12,26,4,2)$. The whole-cell formula of a cell is the Lean-generated canonical CNF at the cell's representative mask, and since every graph of the cell satisfies its formula (the normalization obligation of Section~\ref{sec:strata}), an UNSAT proof of the formula excludes the cell. These five formulas, two for H3 and three for H5, are the inputs of the LRAT consumer of Section~\ref{sec:strata}. The two H3 formulas (29,500 variables and 1,328,183 clauses each) carry checked LRAT proofs, and H3 is excluded a second time by paper reductions with exhaustively replayed finite searches, so its SAT rows are cross-checks.

\paragraph{H3 and H5: the cube route, and how H5 was closed.} The cube consumer of Section~\ref{sec:strata} takes a different input set, seven base formulas: the four H3 scout formulas \lean{orderFortyNineGeneratedThreeHighDistOneB1ScoutCnf}, \lean{orderFortyNineGeneratedThreeHighDistOneC1ScoutCnf}, \lean{orderFortyNineGeneratedThreeHighDistOneC2ScoutCnf} and \lean{orderFortyNineGeneratedThreeHighDistTwoScoutCnf}, and the three H5 cell formulas \lean{orderFortyNineGeneratedVariableHighSatCnf} at \lean{orderFortyNineFiveHighT0Masks}, \lean{orderFortyNineFiveHighT1Masks} and \lean{orderFortyNineFiveHighT2Masks}. Each base formula is covered by two checked cover formulas and a $7\times 8$ grid of 56 positive cubes, a cover in the sense that every graph of the corresponding cell satisfies one of the cubes or one of the cover formulas. Lean proves the accounting, $4\times 56+3\times 56=392$ positive cubes plus $8+6=14$ negative covers, 406 bounded jobs in all (\lean{orderFortyNineSmallHigh_positiveCube_job_count}), and their exhaustive composition (\lean{orderFortyNineSmallHigh_unsat_of_checkedCubeGrid}), so a host runs bounded jobs instead of trusting a solver-side cube list. The five whole-cell formulas and the seven base formulas are alternative inputs to the same conclusion and are never added together. The 2026-08 campaign fired that grid, and the frozen H5 inventory of 2026-09-10 (\repofile{research/problems/erdos-85-wip-01/PHASE_B_H5_H7_INVENTORY_20260910.md}{PHASE_B_H5_H7_INVENTORY_20260910.md}) records 58 roots per H5 cell, 15 of them with direct 2026-08 certificate objects and 43 remaining; the $3\times 43=129$ remaining roots are the H5 rows of the census index, beside the 2 H3 rows, which are the two whole-cell formulas. H5 itself was closed on 2026-09-10 by a different route: reviewed mathematical reductions of the three cells with independently checked finite computations (T0: 1,665 heavy classes reduced to 13 empty negatives plus a separate counting contradiction, review 2037; T1: 249 heavy classes reduced to none, review 2032; T2: 13 heavy classes, 12 excluded by the reviewed census chain and the 44-core with its 92 branches each assigned once to a reviewed exclusion, reviews 2062 and 2063), joined by an outer ledger (\repofile{research/problems/erdos-85-wip-01/q7_h5_closure_ledger/}{q7_h5_closure_ledger/}) whose checker reconstructs all 49 masks of each cell against the Lean canonical arrays; the outer review 2065 accepted the ledger (PASS, 2026-09-10) with an independent checker and pinned result hashes. This is paper-and-computation evidence under the established graph-reduction premises: the three Boolean-exclusion premises of \lean{orderFortyNineStratumExcluded_five_of_booleanExclusions} are undischarged in Lean, no SAT run and no certificate replay were part of the closure, and the 129 root cubes above were not consumed by it. The H5 rows of the census index are therefore the SAT-side inventory of the same three cells, not the evidence for their exclusion.

\paragraph{H7: cells and certificates.} With seven high vertices the size-three supports form a linear triple system on seven points, and \lean{orderFortyNine_highIncidence_profile_of_seven_high} bounds their number $t$ by 7, giving the cells $t=0,\dots,7$ of Table~\ref{tab:inputs}. A kernel-checked census (\repofile{proofs/Proofs/Erdos85OrderFortyNineSevenHighCanonicalCensus.lean}{Erdos85OrderFortyNineSevenHighCanonicalCensus.lean}) shows that every such system is permutation-equivalent to one of fourteen canonical representatives, distributed $1,1,2,3,3,2,1,1$ over $t=0,\dots,7$, and Lean proves the graph cover (\lean{sevenHighCanonicalGraphCover_all}): every $C_4$-free graph on 49 vertices with minimum degree 7, exactly seven high vertices and $t$ low vertices of support 3 satisfies the Boolean constraints of one representative of block count $t$. The thirteen representatives with $t\ge 1$ are excluded by packed LRAT certificates checked inside Lean. Each of the thirteen modules, \repofile{proofs/Proofs/Erdos85OrderFortyNineSevenHighT1Rep0Certificate.lean}{Erdos85OrderFortyNineSevenHighT1Rep0Certificate.lean} through \repofile{proofs/Proofs/Erdos85OrderFortyNineSevenHighT7Rep0Certificate.lean}{Erdos85OrderFortyNineSevenHighT7Rep0Certificate.lean}, reads its certificate (an \verb|include_str| of the packed file) from an absolute path on the project's artifact volume and proves the LRAT check of the formula \lean{orderFortyNineGeneratedH7SatCnf} at the representative's masks (1,329,041 clauses each) by \texttt{native\_decide}. It then derives the representative's exclusion through \lean{sevenHighCanonicalRepresentativeExcluded_of_lrat}. The aggregate \lean{orderFortyNineStratumExcluded_seven_of_t0} (\repofile{proofs/Proofs/Erdos85OrderFortyNineSevenHighCertificates.lean}{Erdos85OrderFortyNineSevenHighCertificates.lean}) reduces the whole stratum to the single $t=0$ representative. The trust profile of these thirteen proofs differs from that of the statements in Table~\ref{tab:axioms} and is stated here rather than there. The checks depend on \texttt{Lean.ofReduceBool} in addition to the standard axioms. The modules were compiled at the 2026-08-15 commits and were not rebuilt in the cold audit of 2026-09-27, so their compile receipts are the commit records and the production manifests, which record for each instance an external drat-trim verification, an external lrat-check verification and a Lean replay of the compact LRAT at production time. The packed certificates reside on the project's artifact volume, which is not published, so the thirteen modules cannot be rebuilt from the public source tree alone (Section~\ref{sec:artifacts}).

\paragraph{H7: the $t=0$ cell.} What remains for the stratum is the $t=0$ cell, whose Lean interface is the uninstantiated capstone \lean{orderFortyNineStratumExcluded_seven_of_emptyCubeEvidenceVectors} of Section~\ref{sec:strata}, with its four evidence vectors of lengths 19, 15, 7 and 2. In that cell the 42 low vertices split into seven of support 0 (\emph{empty}), fourteen of support 1 (\emph{singleton}) and twenty-one of support 2, and the closure works from the induced graph on the seven empty vertices, which is $C_4$-free with maximum degree 3 and falls into 43 isomorphism classes with $a=6$, 7, 8 or 9 edges (the empty-block classification). The universal singleton-capacity argument excludes 15 of the 43 classes without any triangle-count or spectral assumption: the fourteen singleton vertices force at least $\max(0,35-4a)$ common-neighbour pairs outside the empty block, such pairs must avoid the existing empty common neighbours, and a vertex of degree $d_E$ into the empty block can carry at most $7-2d_E$ of them; a vertex-subset upper bound contradicts the lower bound in 12 classes at $a=6$ and 3 at $a=7$, leaving 7, 12, 7 and 2 classes at $a=6,7,8,9$, the 28 structural roots. The generic capacity inequality (\repofile{proofs/Proofs/Erdos85VertexSubsetEdgeCapacity.lean}{Erdos85VertexSubsetEdgeCapacity.lean}), the exterior-pair capacity $\deg_X v+2\,n_E(v)\le 7$ (\repofile{proofs/Proofs/Erdos85OrderFortyNineSevenHighT0ExteriorPairCapacity.lean}{Erdos85OrderFortyNineSevenHighT0ExteriorPairCapacity.lean}) and the total-edge inequality $35\le 4a+|X|$ are Lean theorems with standard axioms; the 43-class enumeration, the representative certificates and their isomorphism transfer are not Lean-certified. The closure ledger of 2026-09-15 (\repofile{research/problems/erdos-85-wip-01/H7_CLOSURE_20260915.md}{H7_CLOSURE_20260915.md}, with its review archive \repofile{research/problems/erdos-85-wip-01/h7-closure-ledger-20260915/}{h7-closure-ledger-20260915/}) then joins each of the 28 roots to an accepted structural exclusion (source cover, high-vertex assignment and quotient cover for the $a=6$ family; source enumeration, with complement completion where the ledger required it (three roots), for eleven of the twelve $a=7$ roots and a cycle-exclusion chain for the twelfth; singleton normalization for $a=8$; maximum-high normalization with a twin and a crossed census for $a=9$), with no residual root. The two capped frontiers were closed by exact partition of their host leaves; for the last root, F14, the 2,278,608 host leaves partition as $1{,}757{,}882+75{,}027+445{,}699$, and the final 445,699 certificates were independently audited and reproduced by a third, independently written implementation with zero mismatches. The premises the ledger cites as established are the degree-7/8 setting with independent high vertices (reviews 1573 and 1574; both are also Lean theorems, Section~\ref{sec:strata}) and the singleton--empty sums and distinct high colours used by the $a=6,7$ source arguments (review 2091). This closes the $t=0$ cell at the level of reviewed prose plus reviewed computation, with no Lean statement for the closure itself: the completeness of the $a=7$ source enumeration rests on an accepted audit of the enumerator code, not an independent exhaustive replay, and the Lean capstone is uninstantiated. Together with the thirteen certificates, the H7 stratum is excluded at the paper-and-computation level with a Lean-checked reduction to its last cell.

\paragraph{H1: the instances.} An H1 graph has a unique high vertex whose eight neighbours induce a perfect matching; every other vertex has exactly one neighbour among those eight, and the remaining 40 vertices induce a 6-regular $C_4$-free graph with eight attachment groups of five. The Lean reduction \lean{orderFortyNineStratumExcluded_one_of_pureFamilies} works through five ordered family profiles, and the frozen candidate list distributes its 1,257 tags over them as 283, 346, 388, 198 and 42. Each candidate carries the 24 table values from which the Lean-emitted binary regenerates its canonical DIMACS (\texttt{v2cnf emit <profile> <table-file>}) and re-checks it (\texttt{v2cnf check}); Section~\ref{sec:trust} says what that pins and what it does not.

\paragraph{H1: the two-solver census.} The frozen Phase B index has 1,416 cases: 96 H1 historical roots with 2026-08 drat-trim-verified certificates, 1,161 H1 residual roots with no accepted certificate object, and the 2 H3, 129 H5 and 28 H7 rows discussed above; H1 is the first two classes, 1,257 instances. Every residual root was solved verdict-only, with no proof logging, by Kissat 4.0.4 and then, on a Kissat UNSAT, by CaDiCaL 3.0.1, each under a declared per-solver cap; a row counted only when both solvers returned UNSAT. Rows that reached a cap were rerun at a longer cap (4, then 12, then 24 hours, the reviewed maximum). Table~\ref{tab:passes} lists the passes; the outcomes sum to 1,160 whole-instance UNSAT rows. Six spot reclaims restarted their rows from scratch. Actual cloud cost for 2026-09-21 to 09-27 was \$235. Every one of these rows was later re-solved with proof logging and certificate-checked (Section~\ref{sec:cost}).

\begin{table}[htbp]
\centering\small
\caption{H1 verdict-only census passes (2026-09-21 to 09-27). Caps are per solver.}
\label{tab:passes}
\begin{tabularx}{\textwidth}{@{}>{\raggedright\arraybackslash}p{.14\textwidth}>{\raggedleft\arraybackslash}p{.07\textwidth}>{\raggedright\arraybackslash}p{.21\textwidth}>{\raggedright\arraybackslash}p{.20\textwidth}>{\raggedright\arraybackslash}X@{}}
\toprule
Pass & Rows & Cap & Hosts & Outcome \\
\midrule
Pilot (09-21/22) & 24 & 4 h & Mac, 4 workers & 24 UNSAT \\
1 (09-21/23) & 1,137 & 4 h & 4 spot c7g.16xlarge & 876 UNSAT, 256 cap hits, 1 infrastructure error, 4 killed in flight \\
2 (09-23/25) & 261 & 12 h & spot, then on-demand & 242 UNSAT, 19 cap hits \\
3 (09-25/26) & 17 & 24 h & 1 on-demand c7g.8xlarge & 16 UNSAT, 1 cap hit \\
3b (09-25/27) & 2 & 24 h & 1 on-demand c7g.xlarge & 2 UNSAT \\
4 (09-26/27) & 1 & 1 h probe, 24 h leaves & Mac, 24 workers & UNSAT through 36 cubes \\
\bottomrule
\end{tabularx}
\end{table}

The one row that defeated every whole-instance cap, \lean{h1_81494a6ef36d3ec9}, had also failed the 1-hour and 4-hour caps of the 2026-08 fleets. Its base CNF was regenerated by the reviewed materializer (42,160 variables, 613,228 clauses, sha256 \texttt{860d8af2\ldots}). The adaptive cube-and-conquer procedure pre-splits five levels by unit-propagation lookahead, probes each cube with Kissat for one hour, confirms each Kissat-UNSAT cube with CaDiCaL under a 24-hour cap, and splits any unresolved cube on the best lookahead variable. The result is a tree of 71 nodes, 35 splits (31 initial, 4 after a failed probe), 36 leaves and maximum depth 8, every leaf UNSAT under both solvers, in 9.7 Kissat and 6.2 CaDiCaL core-hours and 4 hours 50 minutes of wall clock on 24 cores; 9 leaves fell in under 10 seconds, 29 in under 15 minutes, and the hardest needed 0.87 hours of Kissat and 0.96 hours of CaDiCaL. Because each split replaces a cube by two cubes that partition it, the leaves partition the base formula. A separate checker (\repofile{research/problems/erdos-85-wip-01/phase_b_h1_verdict_cloud_20260921/cube_tree_check.py}{cube_tree_check.py}, whose output is \repofile{research/problems/erdos-85-wip-01/phase_b_h1_census_20260927/cube-h1_81494a6ef36d3ec9/tree-check.json}{tree-check.json}) re-derives the tree, every cube's bytes and every verdict from the logs and exits cleanly. An earlier fixed five-variable split on top-occurrence variables refuted 27 of 32 cubes within 3 seconds and was abandoned as uninformative; it is kept only as a datum for Section~\ref{sec:cube}.

The residual-root count comes from exact tag and CNF joins of the solver receipts by the reviewed summarizer over 1,413 cloud run directories and the pilot, not from a sum of overlapping inventory counts. All 1,161 residual roots are UNSAT under two solvers; there were no SAT candidates and no solver disagreements. A second framing, the capacity-grid snapshot of 13,351 slots with 1,288 certificate gaps, is an overlapping decomposition of the same evidence: 1,158 gap slots are residual roots above, 96 are the historical overlay, and 34 lie outside the frozen Phase B source and were solved through their own reviewed route on 2026-09-27/28 (34 of 34 UNSAT under both solvers, Kissat 0.3 to 1.7 hours per row, no cap hits). The three residual roots that are not gap slots are the inventory's three ``historical object conflicts'', which we read as follows: the capacity snapshot already lists a certificate object for each, so the gap auditor does not count them as gaps, but the frozen candidate set retained them and the census re-solved them. The reviewed gap auditor reports 1,191 slots with fresh two-solver UNSAT verdicts, 96 with inherited 2026-08 certificate verification, and one slot, the cube-partitioned row, which it cannot read and lists as UNKNOWN; by its own definition only fresh verdicts close a slot, so it prints 97 ``open'' slots. Every slot has evidence. Since October 2026 every residual root and every historical root also has a checked certificate (Section~\ref{sec:cost}); the 34 slots outside the frozen source have verdicts only.

\subsection{Evidence levels and formal obligations}\label{sec:evidence}

The evidence classes in Table~\ref{tab:evidence} are not a partition: the 96 historical roots, the certificate-object inventory and the Phase B root set use different selections of the same capacity tags. No row is promoted from a solver verdict or an archived proof object to a kernel-checked theorem.

\begin{table}[htbp]
\centering\small
\caption{Evidence level of each input to Result A, and the limit of what each establishes for this paper.}
\label{tab:evidence}
\begin{tabularx}{\textwidth}{@{}>{\raggedright\arraybackslash}p{.17\textwidth}>{\raggedright\arraybackslash}X>{\raggedright\arraybackslash}p{.29\textwidth}@{}}
\toprule
Class & Current evidence & Limit for this paper \\
\midrule
Order-48 and order-49 lower witnesses & Finite graphs in \repofiletab{proofs/Proofs/Erdos85FiniteDropWitnesses.lean}{Erdos85FiniteDropWitnesses.lean}, elaborated through \texttt{native\_decide}; cold-rebuild audit lists three \texttt{native\_decide} axioms per witness; codegree checks corroborate both graphs, and the edge count (168) corroborates the order-48 graph & Establish the lower sides only; not standard-axiom-only. \\
H3 & Checked LRAT proofs of both whole-cell formulas (\leantab{orderFortyNineGeneratedCanonicalSatCnf} at \leantab{threeHighRepresentativeMasks}); an independent paper-and-computation exclusion & Checked proof objects, not a current kernel replay; the generated formal aggregate is absent. \\
H5 & Reviewed reductions of the three cells with independently checked finite computations (reviews 2037, 2032, 2062 and 2063; outer ledger review 2065 PASS, 2026-09-10); no SAT run or certificate replay; the Lean interface \leantab{orderFortyNineStratumExcluded_five_of_booleanExclusions} takes three Boolean-exclusion premises & Paper-and-computation closure; the three Boolean-exclusion premises are undischarged in Lean. \\
H7, cells $t\ge 1$ & Thirteen packed LRAT certificates checked inside Lean by \texttt{native\_decide} (2026-08-15 manifests with drat-trim, lrat-check and Lean replay at production time); \leantab{orderFortyNineStratumExcluded_seven_of_t0} reduces the stratum to the $t=0$ representative & Proofs depend on \texttt{Lean.ofReduceBool}; not covered by the cold rebuild of Table~\ref{tab:axioms}; certificates on the unpublished artifact volume. \\
H7, cell $t=0$ & All 28 surviving structural roots covered after 15 singleton-capacity exclusions (closure ledger of 2026-09-15) & Paper-and-computation closure; the Lean evidence-vector capstone is uninstantiated. \\
H1 historical overlay & 96 reviewed historical rows with archived drat-trim verification; re-checked 2026-10 (archived DRAT to LRAT by drat-trim, then 50 accepted by the Lean standard library's compiled \texttt{LRAT.check} and 46 by \texttt{cake\_lpr}; none rejected) & Compiled or CakeML-verified checking, not a Lean kernel replay. \\
H1 certificate bank & 12,019 historical ready certificate inputs in the replay sizing model & Object and metadata availability do not establish a completed kernel replay. \\
H1 residual roots & 1,161 Phase B roots with no accepted certificate object: 1,160 UNSAT from Kissat 4.0.4 and then CaDiCaL 3.0.1 under declared caps over the pilot plus four cloud passes; one UNSAT on all 36 leaves of an exact cube partition under both solvers; receipts re-derived by the reviewed summarizer and cube checker; certified 2026-10: every root re-solved with proof logging and its LRAT proof accepted by \texttt{cake\_lpr} (1,160 whole instances; the 36 leaves, composed by a Lean lemma) & Certificates checked outside Lean by a CakeML-verified checker and then discarded (Section~\ref{sec:cost}); the H1 semantic bridge remains open. \\
H1 capacity-grid slots & 1,288 gap slots: 1,158 residual roots, 96 historical, 34 outside the frozen source and solved by their own route; auditor reports 1,191 fresh, 96 inherited, 1 UNKNOWN (the cube row) & The auditor counts only fresh verdicts, so it prints 97 open slots; every slot has evidence, and the 1,254 residual-root and historical slots also have checked certificates (2026-10). \\
\bottomrule
\end{tabularx}
\end{table}

The complete set of two-solver UNSAT verdicts supports Result A computationally and leaves the formal nonexistence theorem unproved. The mathematical dependency chain does not depend on how the computation was scheduled: the exact-value theorem still takes \hnoFortyNine, and no completed-certificate count by itself supplies it. A future formal closure would require the exact cover and semantic bridges, certificate checks admitted into Lean for every necessary row (the H1 checks of Section~\ref{sec:cost} run outside Lean), the generated aggregate, and a cold build with literal axiom audits of the two public endpoints. None of these steps follows from UNSAT verdicts or external certificate checks alone.

\subsection{Where the residual trust sits}\label{sec:trust}

The weakest link in any SAT-based nonexistence result is the encoding, not the solver. Three facts bound that risk here. The H1 instance generator is itself a Lean program, \repofile{proofs/Proofs/Erdos85OneHighV2CnfEmit.lean}{Erdos85OneHighV2CnfEmit.lean}, compiled to the native \lean{v2cnf} binary (sha256 \texttt{4bd9604c\ldots}) that every run pins; the emitted DIMACS is regenerated and compared by the binary's own check mode before any solver starts; and the census pipeline compares each emitted CNF against the hash recorded by the independent 2026-08 producer, refusing any mismatch: of the 1,412 preparation receipts in the four cloud passes, 1,322 matched a recorded producer hash byte for byte and 90 were inputs for which the index carried no producer hash. What the encoding does not yet have is a Lean-elaborated semantic bridge for H1: for H3 and H5 the checked formulas are the Lean terms \lean{orderFortyNineGeneratedCanonicalSatCnf} applied to the representative masks, whereas for H1 the row-to-stratum assembly is an open formal obligation. An encoding error would be systematic and invisible to solving; the order-64 case study in Appendix~\ref{app:collab} records one such error that review caught. After the encoding came the solvers; for H1 they are no longer in the trust base. Every H1 instance has an LRAT proof accepted by \texttt{cake\_lpr}, whose soundness is proved in HOL4 and carried to machine code by the verified CakeML compiler \citep{tan2021cakelpr}, so the H1 trust now sits in that checker's specification and compiled binary and in the identity of the checked CNF, not in Kissat or CaDiCaL. Then come the thirteen H7 certificate modules, whose \texttt{native\_decide} checks depend on \texttt{Lean.ofReduceBool}, which read their certificates from the unpublished artifact volume, and which the cold audit of 2026-09-27 did not rebuild; the H7 source enumeration, which rests on an audited program rather than an independent replay; the uninstantiated H7 Lean capstone for the $t=0$ cell; the three undischarged Boolean-exclusion premises of the H5 interface; and the 96 historical H1 certificates, which were re-checked in October 2026 by compiled or CakeML-verified checkers rather than by a Lean kernel replay.

\subsection{Certifying H1 (October 2026)}\label{sec:cost}

Every H1 instance now carries a checked certificate. The method is check-then-discard, the pattern of the largest SAT-based results \citep{heule2016pythagorean,heule2018schur}: for each of the 1,160 whole-instance rows the CNF is regenerated by the pinned \lean{v2cnf} emitter inside the pinned image and its sha256 is required to equal the census hash; CaDiCaL 3.0.1 re-solves it with binary LRAT proof logging; and the proof is streamed through a hashing relay into \texttt{cake\_lpr} \citep{tan2021cakelpr} as it is produced. A row counts as certified only when CaDiCaL exits with \texttt{s UNSATISFIABLE} and \texttt{cake\_lpr} prints \texttt{s VERIFIED UNSAT}. Proofs are not stored; each receipt keeps the proof's sha256 and length, so anyone who re-runs the identical solver binary on the same CNF can confirm that they regenerated, and then check, the same proof. All 1,160 rows were certified: 5.64 TB of binary LRAT, the largest proof 36.2 GB, 3,435 solver CPU-hours and 207 checker CPU-hours, the check running concurrently with the solve. No proof was rejected and no solve returned SAT. Three rows certified twice, on different Graviton instance types with the same binary, gave byte-identical proofs of 27 to 28.6 GB; a macOS build of the same CaDiCaL release gives different, equally valid proofs, so reproduction pins the exact binary.

The hardest row, \lean{h1_81494a6ef36d3ec9}, was certified through its cube tree: all 36 leaves were re-solved with proof logging and accepted by \texttt{cake\_lpr} (30 also by the Lean standard library's compiled \texttt{LRAT.check}), and the Lean lemma \lean{cnf_unsat_of_h1Cube81494a} composes them, with the tree and table as Lean terms, the cover by kernel \texttt{decide}, and the base formula the Lean-generated \lean{oneHighFamilyV2SatCnf}; its \texttt{\#print axioms} lists \texttt{propext} and \texttt{Quot.sound} only, so the leaf checks are the only external step for that root. The 96 historical certificates were re-checked from their archived DRAT files: drat-trim converted each to LRAT against the canonical CNF, and all 96 were accepted (50 by \texttt{LRAT.check}, 46 by \texttt{cake\_lpr}). The cloud cost was about \$200 under an operator ceiling of \$300, on AWS Graviton spot and on-demand hosts, with the three longest rows finished on a local machine. Seven spot reclaims and three infrastructure faults cost only reruns: no receipt was certified from incomplete evidence, and every affected row was re-run to a certified receipt. The per-row receipts, the pinned tool identities and the reproduction recipe are listed in Section~\ref{sec:artifacts}.

\subsection{Why the whole-instance route sufficed}\label{sec:cube}

Before the run, our projection assumed that whole-instance certificates would be too large to check: the Lean standard library's \texttt{LRAT.check} holds the proof in memory, and on the pilot it needed about ten times the binary proof size in RAM (23.5 GB for a 2.3 GB proof), which made the multi-gigabyte proofs of the long rows impractical and pointed to cube-and-conquer \citep{heule2011cube}, in which each cube yields a small certificate and a short Lean lemma composes the leaves. \texttt{cake\_lpr} removed that constraint. It reads the proof as a stream with a bounded heap: 4 GB sufficed for most rows and 16 GB for the longest, independent of proof length, and it checked at about 20 MB of binary LRAT per second, using about 6\% of the solver's CPU time. Whole-instance proofs of up to 36 GB were therefore checked directly, and the cube route was needed only for the one row that no whole-instance solve had finished. Proof size tracked solver work at about 1.6 GB per CaDiCaL hour for whole instances, against about 6.3 GB per hour for the pilot's cube leaves, so splitting would also have multiplied the bytes to check.

\section{Theorem B: a one-proposition reduction of the infinite problem}\label{sec:theoremB}

\subsection{Statement and reduction chain}\label{sec:chain}

\begin{hypAREG}
For every $k\ge 3$ there is no simple $C_4$-free $2^k$-regular graph on $2^k\cdot 2^k=4^k$ vertices. (Lean: \lean{BinarySquareRegularExclusion}.)
\end{hypAREG}

\begin{theoremB}
$\AREG\Rightarrow\neg\,\Qn$. Lean proves this as \lean{not_erdos85Question_of_binarySquareRegularExclusion} from the standard axioms alone (Table~\ref{tab:axioms}).
\end{theoremB}

The reduction is complete in the following sense. The negation of Erdős 85 is equivalent to the existence of arbitrarily late strict drops (\lean{erdos85Negation_iff_not_question}). A drop is produced by a pair of adjacent orders with an existence half and a nonexistence half: a $q$-regular $C_4$-free graph on $q^2-1$ vertices together with the nonexistence of a $C_4$-free graph of minimum degree $q$ on $q^2$ vertices gives $f(q^2)<f(q^2-1)$ (\lean{PlaneOrderDropWitness.strict_drop}), and cofinally many such pairs refute the problem (\lean{not_erdos85Question_of_cofinalPlaneOrderDropFamily}). For $q=2^k$, $k\ge 3$, the existence half is uniform: deleting the absolute nucleus from the even-characteristic polarity graph gives the witness (\lean{Polarity.c4FreeMinDegreeWitness_even_delete_absolute_nucleus}), and the resulting orders are cofinal (\lean{cofinalEvenFieldSquareExclusion_of_binary}). On the nonexistence half the normalized square-order reduction is an equivalence (\lean{squareOrderTightCoreExists_iff_witness}, \lean{binarySquareOrderTightCoreExclusion_iff}), and parity forces every tight core to be regular (\lean{squareOrder_regular_of_even}). Lean then proves directly that $\AREG$ implies the normalized tight-core exclusion (\lean{binarySquareOrderTightCoreExclusion_of_regularExclusion}) and hence $\neg\,\Qn$. Every arrow from $\AREG$ to the headline is a dependency of Theorem B and so compiled in the cold rebuild of Table~\ref{tab:axioms}; $\AREG$ itself remains an unproved hypothesis. Theorem B is a complete checked reduction, not a proof of $\AREG$, and $\AREG$ is not shorthand for several hidden assumptions.

\subsection{The residue beneath A-REG}\label{sec:residue}

The defect operator makes the residue precise. For a hypothetical $q$-regular $C_4$-free graph on $q^2$ vertices with adjacency matrix $A$,
\[
A^2=(q-1)I+J-D,
\]
and $A$ commutes with $D$ (i). The components of the defect graph $D$ have orders $q\,m_c$ with $\sum_c m_c=q$, and every vertex of component $c$ has exactly $m_c$ graph-neighbours inside that component (ii). Unit parts are impossible for even $q$ (iii), and bipartite defect components are impossible when $4\mid q$; the per-component form of that statement discharges the nonbipartite hypothesis for every regular binary-square candidate (iv). The Lean statements behind (i)--(iv) are, in the same order:
\begin{flushleft}\small
(i) \lean{adjMatrix_sq_eq_sub_secondOrderDefect_of_regular}, \lean{adjMatrix_comm_secondOrderDefect_of_regular};\\
(ii) \lean{binarySquare_regular_exists_defectComponent_partition}, \lean{binarySquare_regular_degree_induce_defectComponent_eq_part};\\
(iii) \lean{binarySquare_regular_no_sizeQ_defectComponent_of_even};\\
(iv) \lean{binarySquare_regular_no_bipartite_defectComponent}, \lean{binarySquare_twoPow_regular_no_bipartite_defectComponent}, \lean{binarySquare_twoPow_regular_defectGraph_not_bipartite}.
\end{flushleft}
What remains is exactly $\NONBIP$: all partitions $q=\sum_c m_c$ with every $m_c\ge 2$ and every defect component non-bipartite.

The one-component subcase has a sharp equivalent form.

\begin{nonbipconn}
For binary $q\ge 8$, every loopless $q$-regular $C_4$-free adjacency matrix $A$ of order $q^2$ is singular. Equivalently, its defect graph $D$ is not connected, since $\dim\ker A=\#\mathrm{components}(D)-1$ (\lean{binarySquare_regular_finrank_adj_kernel_eq_card_components_sub_one}).
\end{nonbipconn}

The Lean statement that a proof of this case would discharge is \lean{BinarySquareConnectedNonbipartiteExclusion}. Its connectedness hypothesis is the open part: there is no uniform Lean theorem routing all disconnected or articulated defect graphs into that statement (such a routing exists only in the separate $q=9$ analysis), so the connected case and the mixed non-bipartite partitions are sibling open cases, and connectivity must not be cited as a proved uniform reduction. The formulation includes all graph hypotheses; proving singularity from generic regularity or spectrum alone would not suffice. The current mathematical frontier is the implication from the full symmetric $C_4$-free incidence completion to that singularity statement.

\subsection{Evidence for and against A-REG}\label{sec:areg-evidence}

The evidence is mixed. Even-characteristic polarity graphs supply the cofinal existence half, and the binary-square defect calculus removes unit and bipartite components. Against that, the Lean-checked exact values $f(15)=f(16)=5$ (\lean{minDegreeForC4_fifteen}, \lean{minDegreeForC4_sixteen}, with lower sides from the explicit 4-regular $C_4$-free graphs \lean{fifteenRegular} and \lean{sixteenRegular}) show no drop from order 15 to 16 and refute the $q=4$ analogue of $\AREG$, since \lean{sixteenRegular} is a $C_4$-free 4-regular graph on 16 vertices; and every attempted generic terminal for the remaining non-bipartite completion problem has failed or exposed a weaker hypothesis. The negative map in Appendix~\ref{app:negative} records those failures as scope statements: determinant and spanning-tree controls, real spectral controls, packing and fractional-transversal relaxations, the inverse-sign separation, Pfaffian and Ihara variants, circulant and divisible-design families, Baer-type repairs, and the finite censuses at orders 64, 78 and 80 and in Cayley families of orders 80, 120 and 168. None of them closes $\NONBIP$, and each says exactly which advertised hypotheses were too weak. $\AREG$ is the live mathematical frontier, not a conclusion licensed by the finite data; the plane-order reading is the principal rival, under which special orders may organize both the constructions and the obstructions without the binary regular-exclusion pattern persisting uniformly.

\section{Interpretation: a computational drop, an unproved infinite frontier, and the price of certainty}\label{sec:interpretation}

Result A is computational evidence, not a theorem. The lower sides of both values are explicit witnesses elaborated in Lean through \texttt{native\_decide}, which adds six named trust axioms beyond Lean's standard three; only Theorem B is standard-axiom-only in the cold-rebuild \texttt{\#print axioms} audit, and the thirteen H7 certificate checks add \texttt{Lean.ofReduceBool} outside that audit. The upper side at order 49 is a Lean-proved case split: H3 and the thirteen positive-triple H7 representatives carry checked LRAT proofs, H5 and the last H7 cell are closed by reviewed arguments with independently checked computation, and every H1 instance, settled first by a two-solver census with no open row, now carries an LRAT proof accepted by a formally verified checker. The statement we can make is therefore that we have strong computational evidence that $f(48)=8$ and $f(49)=7$, and therefore that the threshold drops between two adjacent orders. Erdős Problem 85 asks whether $f(n+1)\ge f(n)$ holds for all large $n$. One drop at 48-to-49 is a data point against monotonicity at small order; it is compatible with either answer to the problem as posed, which a finite computation cannot settle, and we make no claim about the problem itself.

We first chose to stop at belief and publish the price of certainty. The price turned out to be lower than projected: with a bounded-memory verified checker the whole H1 stratum was certificate-checked for about \$200 of cloud compute (Section~\ref{sec:cost}), and the solvers left the H1 trust base. What separates Result A from a kernel-checked theorem is no longer the H1 computation but the formal assembly: the H1 semantic bridge, the three H5 Boolean-exclusion premises, the H7 $t=0$ capstone, and admitting the external certificate checks into Lean. The trust boundary is stated in Section~\ref{sec:trust} rather than hidden.

Theorem B identifies what would turn a single drop into an infinite family. It is deliberately stated as a reduction from the one proposition $\AREG$, and $\AREG$ is an unproved hypothesis with a stated rival. Publishing the twin result makes that disagreement useful: the checked reduction states exactly what a proof of the negative answer must supply, the negative map states which plausible shortcuts do not supply it, and Result A provides a concrete calibration point for future theory.

\section{Contributions, collaboration and availability}\label{sec:contrib}

\paragraph{Contributions and collaboration.} Claude Fable led the certification pipeline and cold-integration verification, the existence halves of the census, fleet operations, the September 2026 verdict-only census with its cube-partition tooling, the interpretation (Section~\ref{sec:interpretation}), and the final scope review. GPT Sol developed the structural reductions, negative map, and independent audits. The October 2026 H1 certificate check (Section~\ref{sec:cost}), including the Lean cube-tree composition lemma, was run by Claude Opus. The project's human operator, who set its priorities and scope, and its shared room infrastructure, which supplied coordination, claims, review gates and a durable transcript, are acknowledged; neither is an author. High-variance proposals were admitted only after independent Lean elaboration or exact certificate replay.

\paragraph{Artifacts and receipts.}\label{sec:artifacts}
Every file, directory and Lean module named in this paper is a link to its location in the Lean Genius repository (\url{https://github.com/rjwalters/lean-genius}) on the branch of record, \texttt{erdos85/integration}. The links point at that branch and will be repointed at the release tag when one is cut. Research files live under \repofile{research/problems/erdos-85-wip-01/}{research/problems/erdos-85-wip-01/} and Lean modules under \repofile{proofs/Proofs/}{proofs/Proofs/}; the labels below are paths relative to those two directories. Every number in Sections~\ref{sec:intro} to~\ref{sec:contrib} traces to one of the following files, except that the plane-order census figures of Section~\ref{sec:areg-evidence} trace to the ledger files linked in Appendix~\ref{app:negative}.
\begin{itemize}\raggedright
\item \repofile{research/problems/erdos-85-wip-01/phase_b_h1_census_20260927/}{phase_b_h1_census_20260927/} (\repofile{research/problems/erdos-85-wip-01/phase_b_h1_census_20260927/CENSUS.md}{CENSUS.md}, \repofile{research/problems/erdos-85-wip-01/phase_b_h1_census_20260927/h1-census-table.json}{h1-census-table.json}, \repofile{research/problems/erdos-85-wip-01/phase_b_h1_census_20260927/h1-census-table.tsv}{h1-census-table.tsv}, \repofile{research/problems/erdos-85-wip-01/phase_b_h1_census_20260927/h1-gap1288-audit.json}{h1-gap1288-audit.json}, \repofile{research/problems/erdos-85-wip-01/phase_b_h1_census_20260927/cube-h1_81494a6ef36d3ec9/tree-check.json}{cube-h1_81494a6ef36d3ec9/tree-check.json}): the reviewed summarizer over 1,413 run directories and the pilot (1,160 two-solver UNSAT rows, 96 historical rows, no disagreement), the 1,288-slot gap auditor (1,191 fresh, 96 inherited, one cube slot), the cube tree, every pass, cap and count of Section~\ref{sec:upper}, and the \$235 cloud cost; the open list is empty.
\item \repofile{research/problems/erdos-85-wip-01/h1_cert_full_20261001/receipts/}{h1_cert_full_20261001/receipts/} (\repofile{research/problems/erdos-85-wip-01/h1_cert_full_20261001/receipts/README.md}{README.md}, \repofile{research/problems/erdos-85-wip-01/h1_cert_full_20261001/receipts/h1_cert_census_receipts.tsv}{h1_cert_census_receipts.tsv}, \repofile{research/problems/erdos-85-wip-01/h1_cert_full_20261001/receipts/h1_cert_census_summary.json}{h1_cert_census_summary.json}): the October 2026 certificate check of Section~\ref{sec:cost} (1,160 rows; 5.64 TB; 3,435 and 207 CPU-hours; the proof hashes, the determinism pairs, the 96 historical and 36 pilot-leaf receipts, the pinned tool identities and the incidents).
\item \repofile{research/problems/erdos-85-wip-01/CENSUS_TIMING_20260928.md}{CENSUS_TIMING_20260928.md}: the per-row final-attempt join (2,704 and 2,996 core-hours, the 1.9-hour median, the 121 and 6 tail rows, the 5,700 core-hour total), the 1,412 preparation receipts for the 1,137 cloud-dispatched rows with the 1,322 and 90 input-identity counts, and the projection worksheet of Section~\ref{sec:cube} with each line's arithmetic.
\item \repofile{research/problems/erdos-85-wip-01/CERT_BANK_STATS_20260926.md}{CERT_BANK_STATS_20260926.md}: the certificate-bank statistics (12,102 rows; 1.2, 4.0 and 26.5 GB; 22.6 TB; the GB-per-solver-hour bins; the 984-second, 509 MB replay rate).
\item \repofile{research/problems/erdos-85-wip-01/sat49/H1_REPLAY_SPOT_16_BUDGET_PLAN_20260916.md}{sat49/H1_REPLAY_SPOT_16_BUDGET_PLAN_20260916.md} and \repofile{research/problems/erdos-85-wip-01/sat49/h1_replay_fleet_costs_20260910.json}{sat49/h1_replay_fleet_costs_20260910.json}: the replay cost model (12,019 inputs, 4,831 and 6,651 box-hours, 17.32 days, \$1,276, 6.06 TB, \$182, the memory gate, and the planned wall days per shard count).
\item \repofile{research/problems/erdos-85-wip-01/sat49/H1_V3_SOLVER_TIMING_20260916.md}{sat49/H1_V3_SOLVER_TIMING_20260916.md}: the v3 worker readback (574 rows; 4,011, 4,539 and 7,102 seconds; 22 capped rows; the 10.2-hour interval as a pipeline interval).
\item \repofile{research/problems/erdos-85-wip-01/AXIOM_AUDIT_COLD_20260927/}{AXIOM_AUDIT_COLD_20260927/} (\repofile{research/problems/erdos-85-wip-01/AXIOM_AUDIT_COLD_20260927/AXIOM_AUDIT_COLD_20260927.md}{AXIOM_AUDIT_COLD_20260927.md}, \repofile{research/problems/erdos-85-wip-01/AXIOM_AUDIT_COLD_20260927/axioms.out}{axioms.out}): the cold-rebuild report and the literal \texttt{\#print axioms} output behind Table~\ref{tab:axioms}.
\item \repofile{research/problems/erdos-85-wip-01/STRATA_AND_SMALL_ORDERS_20260928.md}{STRATA_AND_SMALL_ORDERS_20260928.md}: the Lean names of Section~\ref{sec:strata} and of the order-15 and order-16 values, verified against the source on 2026-09-28.
\item \repofile{research/problems/erdos-85-wip-01/H7_POSITIVE_TRIPLE_CELLS_H5_CLOSURE_FORMULA_COUNTS_20260928.md}{H7_POSITIVE_TRIPLE_CELLS_H5_CLOSURE_FORMULA_COUNTS_20260928.md}: the H7 canonical census and certificate inventory (fourteen representatives; the thirteen certificate modules with their 2026-08-15 manifests and trust caveats), the H5 closure summary, and the five-versus-seven formula counts of Sections~\ref{sec:strata} and~\ref{sec:upper}.
\item \repofile{research/problems/erdos-85-wip-01/q7_h5_closure_ledger/README.md}{q7_h5_closure_ledger/README.md}, \repofile{research/problems/erdos-85-wip-01/q7_h5_closure_ledger/REVIEWED_RESULT.md}{q7_h5_closure_ledger/REVIEWED_RESULT.md}, \repofile{research/problems/erdos-85-wip-01/q7_h5_closure_ledger/review2065.json}{q7_h5_closure_ledger/review2065.json} and \repofile{research/problems/erdos-85-wip-01/q7_h5_closure_ledger/reviewer/REVIEW2065.json}{q7_h5_closure_ledger/reviewer/REVIEW2065.json}: the H5 closure ledger (the per-cell reductions with reviews 2037, 2032, 2062 and 2063, the outer review 2065 with its checker and result hashes, and the statement of what the closure did not include).
\item \repofile{research/problems/erdos-85-wip-01/PHASE_B_H1_H3_INVENTORY_20260910.md}{PHASE_B_H1_H3_INVENTORY_20260910.md}, \repofile{research/problems/erdos-85-wip-01/PHASE_B_H5_H7_INVENTORY_20260910.md}{PHASE_B_H5_H7_INVENTORY_20260910.md}, \repofile{research/problems/erdos-85-wip-01/Q7_H1_H3_SQUEEZE_20260910.md}{Q7_H1_H3_SQUEEZE_20260910.md}, \repofile{research/problems/erdos-85-wip-01/Q7_H5_H7_SQUEEZE_20260910.md}{Q7_H5_H7_SQUEEZE_20260910.md} and \repofile{research/problems/erdos-85-wip-01/H7_CLOSURE_20260915.md}{H7_CLOSURE_20260915.md} (with its review archive \repofile{research/problems/erdos-85-wip-01/h7-closure-ledger-20260915/}{h7-closure-ledger-20260915/}): the H1 profile counts and emitter interface, the H3 and H5 profiles, formulas and root counts, the H7 $t=0$ classification, capacity argument and root closures including the cycle-exclusion chain.
\item \repofile{research/problems/erdos-85-wip-01/PAUSE_HANDOFF_20260927.md}{PAUSE_HANDOFF_20260927.md}: the \lean{v2cnf} emitter hash of Section~\ref{sec:trust} and the trust-boundary summary.
\item \repofile{research/problems/erdos-85-wip-01/sat49/verify_boza48_nonisomorphism.py}{sat49/verify_boza48_nonisomorphism.py} with its receipt \repofile{research/problems/erdos-85-wip-01/sat49/BOZA48_NONISOMORPHISM_RECEIPT_20260928.txt}{sat49/BOZA48_NONISOMORPHISM_RECEIPT_20260928.txt}: the order-48 non-isomorphism check over the graph6 inputs in \repofile{research/problems/erdos-85-wip-01/sat49/data/}{sat49/data/}.
\item \repofile{research/problems/erdos-85-wip-01/manuscript/FIRST_DROP_LITERATURE_CHECK.md}{manuscript/FIRST_DROP_LITERATURE_CHECK.md}: the readings of the Boza table and the Afzaly--McKay records, and the $r$-to-$f$ conversion with its dated correction.
\end{itemize}
The review numbers quoted for H5 and for the H7 $t=0$ cell (Section~\ref{sec:upper}, Tables~\ref{tab:inputs} and~\ref{tab:evidence}) are the identifiers recorded in the linked ledger files, which state each review's outcome, checker and pinned hashes; the reviews themselves are entries in the collaboration transcript, published verbatim (redacted at build time) at \url{https://rjwalters.info/rooms/erdos-85}; the ledger statements summarize them. No requester-pays copy of the certificate bank has been published, so no release manifest is listed.

\paragraph{Data and code availability.} Lean sources are the modules linked above, built with Lean 4.31.0 and the repository-pinned mathlib. The cloud controller, node scripts and cube tools (the adaptive splitter, the per-cube verdict and the tree checker) are in \repofile{research/problems/erdos-85-wip-01/phase_b_h1_verdict_cloud_20260921/}{phase_b_h1_verdict_cloud_20260921/}; the certificate-run tooling (per-row emit, solve and stream-check, the node worker, bootstrap and controller) is in \repofile{research/problems/erdos-85-wip-01/h1_cert_full_20261001/}{h1_cert_full_20261001/} and the pilot drivers in \repofile{research/problems/erdos-85-wip-01/h1_cert_pilot_20261001/}{h1_cert_pilot_20261001/}. H1 proofs were checked and discarded, not stored; they are regenerated by re-running the pinned CaDiCaL binary (Linux arm64 sha256 \texttt{fd601b82\ldots}) on the pinned CNF, and a third-party check needs one caution: \texttt{cake\_lpr} exits with status 0 even when a check fails, so success is the line \texttt{s VERIFIED UNSAT} alone. Two collections are not published: the certificate bank (about 6 TB gzipped) and the artifact volume that holds the full solver logs, CNF snapshots, run directories and the thirteen packed LRAT files the H7 certificate modules read. The collaboration transcript is public at \url{https://rjwalters.info/rooms/erdos-85}.

\appendix

\section{The negative map for A-REG}\label{app:negative}

The search did not merely fail to finish several familiar approaches; it produced exact countermodels or reductions showing why they do not close $\NONBIP$. These are scope statements, not impossibility theorems about every future variant: they say which advertised hypotheses were too weak and identify the missing input, a new invariant that uses the simultaneous completion of all rows to one symmetric $q$-regular $C_4$-free $A$. The algebraic and geometric arguments below are reviewed prose unless a named Lean result is identified; their computational checks are evidence for their stated examples, not universal certificates. Where a public ledger file carries an item's record it is linked; the two consolidated records are the campaign's working proof outline (\repofile{research/problems/erdos-85-wip-01/FINAL_PROOF_OUTLINE.md}{FINAL_PROOF_OUTLINE.md}) and its cuts ledger (\repofile{research/problems/erdos-85-wip-01/CUTS_LEDGER_DRAFT.md}{CUTS_LEDGER_DRAFT.md}), and the pointers below are to sections and rows of those two files.

\begin{itemize}
\item \emph{Determinant.} The identity $\det(A)^2=q^4\,\tau(D)$ makes a square spanning-tree count necessary when $D$ is connected. Controls show that this condition alone does not force singularity (proof outline \S A.5.3(i); cuts ledger row 1); stronger integral or local arithmetic tests are not excluded, and later examples fail precisely such tests.
\item \emph{Real spectrum.} Real spectral controls satisfy the scalar conditions recorded in the cuts ledger but do not establish integral adjacency realizability or compatibility with every odd-trace condition. Hoffman-diagonal and odd-power arguments reject particular proposed spectra (proof outline \S A.5.3).
\item \emph{Packing and relaxation.} Generic packing, Hall, code-LP and fractional-transversal statements omit the completion of all rows to one symmetric adjacency matrix; the strongest faithful self-polar partial interface already has an exact half-integral survivor at $q=6$ (cuts ledger).
\item \emph{Inverse-sign separation.} The proposed P8 separation is impossible: exact product-sign identities force both signs among the relevant inverse entries. The false terminal was retracted, not weakened (cuts ledger row 56; Appendix~\ref{app:collab}).
\item \emph{Variants.} Pfaffian and mod-2 valuation, Ihara, exterior-power and Lefschetz variants reduce to the same determinant or spectral data. Coherent-configuration, projective-completion and line-graph thresholds lack the required bridge; the line-graph premise already fails the $q=4$ incidence control. The ten-route audit and its no-survivor verdict are recorded in \repofile{research/problems/erdos-85-wip-01/NONBIP_CONNECTED_LITERATURE_DIVERGENCE.md}{NONBIP_CONNECTED_LITERATURE_DIVERGENCE.md}.
\item \emph{Central operation.} A core-Lean partial common-neighbour reconstruction (\repofile{proofs/Proofs/Erdos85PartialCentralReconstruction.lean}{Erdos85PartialCentralReconstruction.lean}) identifies the exact graph problem in operation form; a separate prose construction, checked at finite-field examples, refutes same-carrier Hamming stability based only on the total central law's failure rate (\repofile{research/problems/erdos-85-wip-01/PARTIAL_CENTRAL_OPERATION_AUDIT.md}{PARTIAL_CENTRAL_OPERATION_AUDIT.md}). It does not refute every additional operation identity.
\item \emph{Circulant and representability families.} The interval circulant defect family with steps $\{q^2/2,\pm 1,\dots,\pm(q/2-1)\}$ is excluded for every binary $q\ge 4$ by a reviewed cyclotomic argument (\repofile{research/problems/erdos-85-wip-01/l6-cyclic-trace/UNIFORM.md}{l6-cyclic-trace/UNIFORM.md}). Specific order-256 defects fail rational congruence to the identity form, witnessed by local Hasse invariants (\repofile{research/problems/erdos-85-wip-01/l6-representability/}{l6-representability/}); the recorded D16 also fails the mod-2 alternating-root test, and the interval $r=4$ example passes the 2-adic test but fails at odd primes. These are family exclusions with no reduction of arbitrary defect graphs to those families.
\item \emph{Design and polarity repairs.} No proper divisible design graph has odd degree $k\ge 3$ with common-neighbour parameters $\lambda_1=0$, $\lambda_2=1$. Reviewed geometric arguments rule out loopless polarities preserving the Baer-deleted incidence structure, and matching-repair and retained-pencil bounds rule out stated support and deletion budgets; none excludes arbitrary global repairs (\repofile{research/problems/erdos-85-wip-01/baer-repair-gate/}{baer-repair-gate/}, \repofile{research/problems/erdos-85-wip-01/BAER_POLARITY_EXTENSION.md}{BAER_POLARITY_EXTENSION.md}, \repofile{research/problems/erdos-85-wip-01/retained-subplane-gate/}{retained-subplane-gate/}).
\item \emph{Alternating square roots.} Explicit regular connected non-bipartite defect graphs pass the full alternating symmetric square-root test over $\mathbb{F}_2$, but the same family fails the determinant-square condition uniformly; the two necessary tests are separated, with no survivor of their conjunction (\repofile{research/problems/erdos-85-wip-01/mod2-root-gate/}{mod2-root-gate/}).
\item \emph{Flag relaxations and enlargements} (\repofile{research/problems/erdos-85-wip-01/CUTS_LEDGER_DRAFT.md}{CUTS_LEDGER_DRAFT.md}, rows 174--180). The stated five-vertex flag relaxation is feasible uniformly for binary $q\ge 16$ after repairing the original witnesses to satisfy $\operatorname{tr}(AD^2)\ge 0$, for the listed 23 identities and eight Gram families (row 175); the aggregate mod-3 census is feasible for every odd exponent $k\ge 5$ (row 174); the projected fixed-shore bound $z\le q-1$ is refuted by a uniform $z=q$ construction (row 176); and the specified enlargement of an even-order polarity graph to degree $q+1$ on $(q+1)^2-3$ vertices requires at least $q^2/4$ old-edge deletions for $q\ge 16$ (row 180), excluding repairs with only linearly many deletions, not all repairs.
\end{itemize}

\paragraph{The 63-to-64 campaign.} At $q=8$ the possible defect-component partitions are $[2,2,2,2]$, $[3,3,2]$, $[4,2,2]$, $[4,4]$, $[5,3]$, $[6,2]$ and $[8]$, and they are not all excluded. The size-two $\mu=3$ sector is closed on regular hypotheses by \lean{orderSixtyFour_regular_sizeTwoEigenline_false}, and the complete negative signed-joint size-two subtree, including the corrected $\mu=-3,(0,5)$ endpoint, by \lean{orderSixtyFour_regular_sizeTwo_signedJoint_false_of_connected}. But $[3,3,2]$, $[4,2,2]$ and $[6,2]$ retain non-bipartite cases; $[4,4]$ and $[5,3]$ have only partial owner-nullity information; the connected $[8]$ case remains a gap; and the eleven $[2,2,2,2]$ assembly targets are external UNSAT verdicts without certificates. Therefore 63-to-64 is not a decided drop, and none of these order-64 enumerations proves $\AREG$ (proof outline \S A.5.2). The $\mu=-3,(0,5)$ endpoint carries exactly six disclosed Owner88 \texttt{native\_decide} checks (\lean{muNegThreeZeroFiveEndpointCallback_false}); the broader \lean{orderSixtyFour_regular_sizeTwo_signedJoint_false_of_connected} additionally inherits the older FullClosure certificate family, and both axiom lists were printed separately. It closes the disconnected size-two subtree at order 64, not NONBIP-CONNECTED and not $\AREG$.

\paragraph{Plane-order censuses.} The order-80 search for a 9-regular $C_4$-free graph, which would supply the $q=9$ existence half, is unresolved: all 20 solver attempts ended UNKNOWN under their caps, the independently checked positive controls returned SAT but neither is an order-80 witness, and the order-78 existence cases were left open (\repofile{research/problems/erdos-85-wip-01/q9_solver_controls/Q9_EXISTENCE_DECISION_20260911.md}{q9_solver_controls/Q9_EXISTENCE_DECISION_20260911.md}). An UNKNOWN is not evidence of nonexistence. The finite-group Cayley census examined 52 groups of order 80, 47 of order 120 and 57 of order 168 with inverse-closed connection sets of degree 9, 11 and 13; every input returned UNSAT, with reviewed encoding and table audits (\repofile{research/problems/erdos-85-wip-01/cayley-census-q11-q13/CAYLEY_CENSUS_Q11_Q13_20260913.md}{cayley-census-q11-q13/CAYLEY_CENSUS_Q11_Q13_20260913.md}). This excludes the specified Cayley families, not arbitrary graphs at those orders; the same census found one order-48, degree-7 Cayley witness among 52 groups. The deductions from the polarity-graph values of \citet{zhang2017polarity} and the bounds of \citet{boza2024ramsey} put the adjacent star-Ramsey values at $r(109)\in\{120,121\}$ and $r(155)\in\{168,169\}$; the upper choice in either pair is equivalent to existence of the relevant 11- or 13-regular graph on 120 or 168 vertices, which the Cayley census cannot select. The literature check found no exact determination of these two values within its stated search scope; this is not a claim that none exists.

\section{Technical lessons from the collaboration}\label{app:collab}

This appendix records the technical lessons of the two-agent campaign, because the process determined which claims reached the record. The collaboration transcript is public (\url{https://rjwalters.info/rooms/erdos-85}); each lesson is also stated with the Lean modules and ledger files that carry it and can be checked against those.

\paragraph{Verification asymmetry.} Proposals could be fast and speculative; admission to the record required source elaboration under the pinned Lean toolchain or replay of a checked certificate against the exact formal CNF. In the inverse-potential episode a root-summed identity announced as $\operatorname{tr}(A^{-1})=q$ was withdrawn five minutes later, when its author re-expanded the all-pairs term and found $q\cdot\mathbf{1}^{\mathsf T}A^{-1}\mathbf{1}=q^2$ instead; no Lean wrapper or outline node had been built, so the correction cost minutes and left no false theorem behind.

\paragraph{The order-64 case study: an 80-owner error and its 88-owner repair.} The $\mu=-3,(0,5)$ endpoint shows why hypothesis fidelity matters more than an UNSAT line. The first encoding silently reused the h114 shore table (eight fixed owners per shore, 80 candidate owners), whereas the actual h305 shore modes contain antipodal offsets as well (twelve fixed edges per shore, an 88-owner universe); the original 80-owner UNSAT result therefore excluded only a strengthened, wrong problem. The mismatch was detected by comparing the formal shore modes with the certificate universe before the endpoint was called proved. Six canonical 88-owner CNFs were then emitted independently and agreed byte for byte; all six were UNSAT in seconds, and their LRAT payloads are checked by the six \lean{h305Owner88*_check} theorems in \repofile{proofs/Proofs/Erdos85MuNegThreeZeroFiveCorrectOwnerCertificate.lean}{Erdos85MuNegThreeZeroFiveCorrectOwnerCertificate.lean}. The expensive work was not search but proving that the graph semantics map to those exact literals and modes, through \lean{h305_crossExteriorSplit_of_profile}, \lean{muNegThreeZeroFiveCorrectFiniteSemantics_false}, \lean{muNegThreeZeroFiveCorrect_graph_false_of_exterior} and \lean{false_of_h305_source_or_transported}; a cold rebuild of the chain printed the exact six Owner88 axioms. Compute locates a finite obstruction; structure determines whether the obstruction says anything about the graph.

\paragraph{Adversarial diversity.} Different agents were useful only when they tested different failure modes. An independent review re-derived P1--P7 and inspected the exceptional root term, forcing the corrected domain $V\setminus(N(y)\cup\{y\})$; the author then proved that the corrected P8 sign separation was itself impossible, and a second agent caught an overbroad sentence in that retraction and supplied the necessary choice from $N_D(y)\setminus N_A(y)$ (the inverse-sign entry of Appendix~\ref{app:negative}). Independent checking also corrected campaign accounting: a manuscript audit described the thirteen order-49 inputs as five one-high cells plus one seven-high cell, the host manifest listed four H3 scouts, three H5 cells and six H7-t0 cubes, and the sentence was retracted within two minutes. The lesson is narrower than ``use multiple models'': the reviewer must inspect an independent invariant (the exact host manifest, a semantic universe, a boundary term) rather than replay the author's narrative.

\paragraph{The structure--compute exchange rate.} The order-49 campaign gives a measured exchange: structural normalization replaces the five whole-cell H3/H5 formulas of the LRAT interface by seven base formulas (four H3 scouts and three H5 cells), each with two checked cover formulas and a $7\times 8$ grid of positive cubes; Lean proves the accounting (392 positive cubes plus fourteen covers) and their exhaustive composition, and the host runs 406 bounded jobs (Section~\ref{sec:upper}). When the campaign was fired, a separate audit found that these grid results did not fit the older interface that expected five whole-cell LRAT proofs; the missing semantic path was formalized as \lean{not_c4FreeMinDegreeWitness_fortyNine_seven_of_smallHighCubeBaseUnsat} (Section~\ref{sec:strata}) before the first verdict was needed. The certificate campaign of Section~\ref{sec:upper} was fired only after a posted manifest of its thirteen input hashes, tool hashes, job order, caps and storage policy had been verified entry by entry: a verdict without a pre-recorded input hash and a semantic bridge to the formal statement is not progress, even if the process printed UNSAT.

\paragraph{Silence is not success.} The false trace identity above looked plausible because no term visibly objected; an explicit re-expansion produced the positive evidence of cancellation. P8 survived a thousand sampled $q=4$ models only because those models were singular and never instantiated the inverse hypothesis; the correct response was UNKNOWN, not support, and a direct block-sum argument then proved P8 impossible in every hypothetical survivor. Operationally, a solver cell counts only when its ledger line contains a verdict, return code, input hash, proof-check status, LRAT hash and upload status; a Lean file counts only after the named target elaborates and its \texttt{\#print axioms} output is inspected. On 25 August the order-64 dependency chain ran for nearly an hour, and three agents waited for the final file check and separately reported its axiom scope rather than inferring success from the absence of errors. Every status word names its scope: the order-64 callback theorem \lean{muNegThreeZeroFiveEndpointCallback_false} carries the foundational axioms plus exactly six disclosed Owner88 \texttt{native\_decide} checks, the broader \lean{orderSixtyFour_regular_sizeTwo_signedJoint_false_of_connected} inherits the older FullClosure family, and both lists were printed rather than the smaller list for the larger theorem. Positive artifacts, not quiet logs, are the unit of trust.

\paragraph{Methods.} \emph{Room protocol}: a persistent shared chat room (transcript at \url{https://rjwalters.info/rooms/erdos-85}); claim, red-team, formalize, relay-merge, independent cold audit; every push to the working branch is merged to the mirror branch by the counterpart. \emph{Cold-audit rule}: as in Section~\ref{sec:lean-defs}, ``verified'' means the file elaborates from source under the pinned toolchain in an independent environment with verbatim \texttt{\#print axioms} on the public theorems; adopted after a stale-cache incident and tightened after an rg-mask incident. \emph{Certificate factory}: the rules of Section~\ref{sec:trust} (exact emitters with input hashes, two-solver portfolios, DRAT and LRAT replay, durable manifests, and a semantic bridge before a SAT result counts as mathematical evidence). \emph{Census tooling}: exhaustive branch-and-bound sweeps over partition and atom spaces with every constraint tied to a named Lean lemma (loads, budgets, balance integrality, the equal-LCM law, oriented-cover kernels).

\bibliographystyle{plainnat}
\bibliography{refs}

\end{document}
